\documentclass[runningheads]{llncs}
\usepackage{xcolor,amsmath,amssymb,booktabs,longtable,array,url}
\usepackage[T1]{fontenc}
\usepackage[breaklinks,hidelinks]{hyperref}
\urlstyle{tt}
\makeatletter
\g@addto@macro\UrlBreaks{\do\a\do\b\do\c\do\d\do\e\do\f\do\g\do\h\do\i\do\j\do\k\do\l\do\m\do\n\do\o\do\p\do\q\do\r\do\s\do\t\do\u\do\v\do\w\do\x\do\y\do\z\do\A\do\B\do\C\do\D\do\E\do\F\do\G\do\H\do\I\do\J\do\K\do\L\do\M\do\N\do\O\do\P\do\Q\do\R\do\S\do\T\do\U\do\V\do\W\do\X\do\Y\do\Z}
\makeatother
\setlength{\emergencystretch}{2em}
\newcommand{\modelpin}{60af4aaefd64048f188fc9710cc953c5356f26f2}
\newcommand{\checked}{\par\noindent\textcolor{green!35!black}{\textbf{Existing checked declaration.}}\ }
\newcommand{\hand}{\par\noindent\textcolor{blue!65!black}{\textbf{Hand proof supplied.}}\ }
\newcommand{\proposed}{\par\noindent\textcolor{violet}{\textbf{Proposed statement.}}\ }
\newcommand{\missing}{\par\noindent\textcolor{red!65!black}{\textbf{Missing mechanization.}}\ }
\newcommand{\decl}[1]{\textnormal{\small\nolinkurl{#1}}}
\newcommand{\evidence}[1]{\par\noindent\textit{Checked reference:} \decl{Singular.Statements.#1}; exact pin and verbatim axiom result in Appendix~\ref{sec:artifact}.\par}
\begin{document}
\title{Singular: A Permanent Registry with Protected Rejection and Explicit Accounting}
\titlerunning{Singular: Protected Rejection and Explicit Accounting}
\author{}
\institute{}
\maketitle
\begin{abstract}
Singular is a permissionless and reusable mint of named, unique and permanent
tokens on Cardano. These five properties organize the design: a name has one
representative, termination cannot be undone, processing uses public data, and
applications govern their own payloads. Cardano KERI motivates selection of
checkpoints by identifier and evidence of whether an identifier was ever
convicted. These application goals are not established by a registry proof alone.
A requester also needs predictable spending and protection from arbitrary
cancellation. Singular separates a permanent key lifecycle from application state and places
pending requests in independent outputs. This paper describes the ruled design:
registration and termination, booking and token provenance, reference-read
rejection, funded timed retraction, and the allocation of request value. It separates
properties checked in Lean from conditional arithmetic and unresolved ledger
correspondence. Rejection requires a current lifecycle mismatch or expiry after
both registration windows. Termination has only a fold exit. The accounting
analysis distinguishes gross tip allocation from profit and identifies the limits
of refund-only judgements. Neither inclusion, resistance to flooding, nor a
positive processor wage follows. Validators and builders are not yet conformant
to the ruled economics; model proofs do not change that standing.
\end{abstract}
\keywords{registry \and extended UTxO \and formal specification \and
protected rejection \and transaction accounting}

\section{Named, unique, permanent; permissionless and reusable}
\label{sec:introduction}
\begin{quote}
Singular is a permissionless and reusable mint of named, unique and permanent
tokens on Cardano.
\end{quote}
This is the adopted definition and the organizing promise of the design.
The tokens are named, unique and permanent; the mint is permissionless and
reusable. Section~\ref{sec:five-properties} gives each meaning, its concrete
story and the evidence that does, or does not yet, support it. The definition
does not turn application demonstrations into Lean theorems or imply that the
ruled economics is already deployed.

\subsection{Why a name matters}
Cardano KERI motivates two questions: how to select checkpoints by KERI
identifier, and how to prove whether an identifier was ever convicted on
Cardano. For that application a conviction is a termination, and it is
permanent. A lookup of an identifier at a registry root $R$ has three
outcomes. An active leaf proves that the identifier is registered and not
convicted at $R$: ``not convicted'' means active. A terminal leaf proves
conviction. An unknown leaf proves only that the identifier is not
registered on Singular; that is a different answer, and it is never
presented as ``not convicted''. The stored Absent leaf belongs to the
extension edges outside the contract of the first release and is not part of
this reading.
\checked A terminal key stays terminal: for a reachable registry whose leaf
at the key is terminal, every request at the key is refused and every accepted
later batch leaves the leaf terminal.
\evidence{termination}
\checked A terminal attestation, once held, stays held and its key's leaf is
terminal after any accepted fold batch. A holder may separately discard an
attestation without changing the terminal leaf; this is not a fold.
\evidence{terminal_attestation_permanent}
\evidence{terminal_discard_changes_no_registry_state}
\hand By induction over the accepted batches from the genesis state, a key
that is terminal at one reachable state is terminal at every later reachable
state. Hence an active leaf at $R$ shows that the key was never terminal at
any earlier state: an active identifier at $R$ was never convicted before
$R$. The argument uses only the checked one-batch statement and the
reachability of $R$.
\missing The lookup is a logical read of the root; verifying a concrete
lookup proof is a correspondence obligation of the validator or verifier. A
lookup is valid against the current root only until the next fold;
issuance and implementation of durable attestations of presence, absence
and termination remain future work. Singular alone neither verifies a duplicity
proof nor
establishes a KERI checkpoint or conviction history, and the consumer's
integration remains open.

\subsection{The reader's problem and the evidence boundary}
A requester wants to know which funds a failed registration returns. An
application wants its key's representative to remain unique while its own state
changes. A processor wants to know what it may consume, what it must pay and
which evidence authorizes each exit. A scientific reviewer needs these promises
separated from their implementation and economic assumptions.

The registry records a lifecycle, not key ownership or application payload.
Only registration of an unknown key and permanent termination of an active key
change it. Applications authorize entry and govern the resulting token in their
own custody. Processing is permissionless within the admission rules: a stranger
cannot simply cancel a compatible registration before expiry. A current
lifecycle mismatch is independently sufficient for rejection, including during
the reclaim window. This distinction matters when two requests race for one key.

The behavioral authority considered here is the Lean source at the exact revision
\begin{center}\small\texttt{\modelpin}\end{center}
using Lean 4.27.0. This is the last ancestor revision changing the model source.
The appendix binds checked claims to fully qualified declarations and their
fresh axiom-check output. A reference to a declaration is a claim only about its
written hypotheses and conclusion, not about its name, a deployed script or a
ledger execution.

Four labels distinguish evidence throughout. \textbf{Existing checked declaration}
means the appendix contains a declaration and its compiled axiom result at the
pin. \textbf{Hand proof supplied} means the argument appears here with its
premises. \textbf{Proposed statement} identifies a conditional design or sizing
criterion without measurements establishing it. \textbf{Missing mechanization}
identifies a promised boundary for which no checked result is supplied. These
labels are independent of implementation conformance.

\section{The five properties and their stories}
\label{sec:five-properties}
The following meanings state the adopted design in plain words. Checked
results concern the specified model; application behavior and a running
demonstration require evidence at their own boundary. Every checked reference
uses the exact revision and axiom line in Appendix~\ref{sec:artifact}.

\subsection{Named}
Each token has a name: registering a name mints exactly one token for it,
and the registry is the only way to mint it. Holding the token is holding
the name. The concrete story is that a registered name has a token which
an application can hold and associate with its own payload.

\checked The original registry's accepted registration transaction mints
one active asset indexed by the request's key, delivers it to the named
destination and records that holding. Its model asset name equals the key.
\evidence{insert_active_transaction_row}
\checked In an economic reachable world, actual singular-tag supply across
application custody and pending terminations is one exactly for an active
key, and zero otherwise.
\evidence{world_supply_sync}
\missing These results are about modeled identities and transitions.
``Only way to mint'' on Cardano requires correspondence of the actual minting
policy and name encoding. Holding the name's representative is not a proof
of a legal right, personal identity or entitlement to that name. The
application decides what its payload means.

\subsection{Unique}
A name can be taken only with a proof that it is absent, so two people cannot
hold separate genuine tokens for the same name in one registry. Bob tries
to register Alice's name after she has registered it; his registration is
refused.

\checked After an accepted original registration, another registration of
the same key with valid approval is refused by the modeled builder with
\texttt{key-exists}. The same declaration records exactly one active holding.
\evidence{insert_active_transaction_row}
\checked The economic supply invariant retains the bound of one while the
token moves from application custody into a pending termination.
\evidence{world_supply_sync}
\missing A concrete absence proof's verification by the validator is a
correspondence obligation. ``Absent'' here means an unknown, untaken key,
not the historical stored Absent leaf. The refusal concerns taking an
already active name; booking a competing request does not reserve it, and
the model does not adjudicate who ought to have registered it first.

\subsection{Permanent}
Terminating a name burns its token and records it as terminal; it can never
be issued again. Alice terminates her name and then tries to register it
again. The attempt is refused. For KERI, termination is intended to be
conviction, so a convicted identifier must never come back in that registry.

\checked For a terminal key in a reachable original registry, every request
at that key is refused; every accepted later batch keeps it terminal.
\evidence{termination}
\checked A genuine pending termination has an accepted economic fold that
makes its key terminal under the theorem's explicit funding and willingness
premises. World supply is then zero for that key.
\evidence{booked_termination_foldable}
\evidence{world_supply_sync}
\missing The application decides when release from custody is authorized.
A KERI conviction needs its own duplicity-proof and integration evidence.
Permanence within this registry does not identify a new registry with the old
one. Economically, the termination donates everything held except the tip
to the treasury; it returns no deposit to Alice.
\evidence{fold_routes_every_amount}

\subsection{Permissionless}
Anyone creates a registry, anyone books a name, anyone settles pending
requests, all from public chain data. The creator has no power over the
names. These are the design's participation promises, subject to application
authorization, funding and the admission rules. The concrete story is Bob,
who did not create the registry, settling the pending requests, Alice's included.

\checked The original registry builder is equal to a construction from its
public view: the registry and pending request, including the request's datum.
An accepted original fold requires no signer, and changing approval signatures
does not change that step or its built transaction.
\evidence{fold_inputs_public}
\evidence{fold_requires_no_signer}
\missing These are original-builder statements, not a public-data theorem
for the economic transaction builder or an execution of Bob's story. Registry
creation, obtaining all chain data, funding and submitting the actual batch
require off-chain and validator evidence. The story is a demonstration of
that boundary, not a model theorem or an inclusion guarantee. Creator
privilege over application payloads is likewise not decided by these results.

\subsection{Reusable}
The same mint serves any application: open-datum today, a closed variant in
the demonstration, KERI checkpoints next, with no change to its contract.
Each name's token lives in its owner's application, which decides what the
name points to. This is the design's separation of the name from its payload.

\missing The concrete application demonstration uses a closed flag on the
same release and commands: a payload set at registration cannot be updated,
the application refuses every update, and only termination remains. This
is application-validator and executable behavior outside the checked registry
model. No release or running-demonstration receipt is asserted by this paper.
KERI is the next intended consumer, not a completed integration. Serving
both application kinds within one registry through application tags is also
a future integration claim. The application's release rule and adaptations
to the ruled termination economics remain open.


\section{System context and trust boundaries}
The extended UTxO model attaches validation to outputs consumed by a
transaction~\cite{ExtendedUTXO20}. Singular separates three responsibilities.
The application certifies a registration payload and controls release of the
singular tag. Singular checks the registry lifecycle and the request's admitted
exit. An authenticated map implementation supplies concrete membership and
nonmembership evidence. The ledger supplies input availability, validity bounds,
fees, output minima and atomic transaction acceptance.

\begin{figure}[ht]
\centering
\begin{tabular}{c}
\fbox{\parbox{.85\linewidth}{\centering Application boundary: authorize entry;
maintain application datum; release the singular tag}}\\[4pt]
$\downarrow$\quad certified registration or custody release\\[4pt]
\fbox{\parbox{.85\linewidth}{\centering Singular boundary: book requests;
judge their exits; update the lifecycle and token supply}}\\[4pt]
$\downarrow$\quad spending, references, outputs, mint and donation\\[4pt]
\fbox{\parbox{.85\linewidth}{\centering Ledger boundary: authenticate inputs,
run scripts, enforce balance, fees, time and output minima}}
\end{tabular}
\caption{Responsibilities. An application policy's semantic correctness and
the translation of logical observations to ledger data are external obligations.}
\end{figure}

A state output identifies one registry by its exact state token and carries
its configuration and current root. The logical model retains the map itself;
a concrete validator works with authenticated proofs. A fold consumes and
continues the state output. A rejection reads it by reference. Pending requests
occupy separate outputs; their existence alone reserves no turn on the state
output. Application-local updates can preserve the singular tag without
performing a registry transition.

The intended lifecycle is
\[
 \text{Unknown}\ \xrightarrow{\text{registration fold}}\
 \text{Active}\ \xrightarrow{\text{termination fold}}\ \text{Terminal}.
\]
Terminal is permanent. Stored Absent is a distinct historical encoding, not a
synonym for Unknown; five historical operation encodings are refused by the
active two-operation law. No migration or parameter-update authority is part of
this registry. A newly created registry is a different registry, not a migration
of its pending requests.
\checked For a Terminal key in a reachable original registry, every request
at that key is refused and every accepted later batch keeps its leaf Terminal.
\evidence{termination}

The software contains Aiken state, request and witness validators, Haskell
transaction builders and ledger/provider interfaces, the Lean authority, and
conformance consumers. Their presence is not correspondence evidence. The
existing request validator still uses a state-spending rejection path; the ruled
economic judgement is separate. The public suite retains historical receipts
and publishes affected requirements as unmet. The old driver's transactions and
its proofs must not be substituted for the economic transactions described here.

\section{The ledger boundary and staking}
\checked The modeled policies have a complete per-kind decision. Active and
absent movements, burns included, require a fold. Terminal movements require
a fold unless every movement is a burn. A fold spends an input at the state
script's payment credential, with exactly one unit of this registry's state
token and the fold redeemer.
\evidence{witness_decision_cases}
\checked Under bound policy identifiers, ledger script execution and collision
freedom on the scripts occurring in that transaction, an accepted active or
absent movement also establishes that the state script accepted that spend.
Collision freedom is a named local premise, not a global injective-hash claim.
\evidence{fold_required_for_active_and_absent}

\checked A positive active mint is connected to that same spent state, its
datum, an admitted batch, the validity upper bound and the batch's complete
token mint. The key was unknown before the fold; from a reachable state it
has exactly one active token afterward. Besides the premises above, this
uses the explicit \decl{StateScriptFolds} correspondence premise.
\evidence{active_mint_at_unknown_key}
\missing The premise is not a proof that a compiled state validator implements
the model batch. That correspondence remains an implementation obligation.

\checked All three token policies ignore staking: changing withdrawals,
certificates, their redeemers and any input, reference-input or output stake
credential leaves the decision unchanged when the transaction with those
fields erased is the same.
\evidence{token_policies_read_no_staking}
\missing This result is not a staking theorem for the state or request rules;
those remain in the state/request model and correspondence work. A delivered
output still binds the whole address named by its request, stake part included.

\checked Removing one terminal holding preserves every key's leaf, the root,
configuration, custody, every active count and consistency. Terminal
attestations are zero or many; their number is not the active supply invariant.
A holder's discard is separate from a fold, so attestation permanence through
folds remains true. The terminal fact survives discarding its evidence token.
\evidence{terminal_discard_changes_no_registry_state}

\checked One finite ledger satisfies all named premises. It accepts an active
mint spending a reachable registry state, with one active token after the
admitted batch, and a terminal pure burn without a fold that consumes a real
attestation. The discarded holding exists in a consistent terminal state.
\evidence{ledger_boundary_premises_satisfiable}
\missing The example does not establish terminal issuance through the active
two-edge lifecycle or either transaction's acceptance on Cardano. Terminal
issuance and concrete deposit recovery require their own evidence.

\section{Adversary and cost model}
The actors are request funders, refund recipients, application custodians,
processors, block producers and the treasury. One party may occupy several
roles. A destination need not belong to the requester. Treasury receipts are
not the wealth of an individual treasury voter.

An adversary may create many well-formed requests, select or withhold processing,
observe public pending transactions and, in blocks it produces, order competing
transactions. We assume no particular stake fraction, fair inclusion rule,
minimum service rate or rational participation. Cryptographic correspondence
requires authentic signatures and assets and sound authenticated-map proofs.
Those are boundary assumptions, not consequences of the abstract token counts.

Amounts, times and window lengths are natural numbers; signed net changes use
integers. Monetary units are lovelace. Ledger fees are nonnegative.
For an action $x$, let $c_x$ be its actual ledger fee; a separate $e_x$ may cover
computation and operational expense not paid to the ledger. The latter is not a
transaction output. Let $V(k)$ be a party's valuation of key $k$; it is outside
Singular's monetary accounting. Preferences or side payments involving $V(k)$
can dominate the declared tip.

A request holding $H$ for $\Delta$ time units has simple opportunity cost
$rH\Delta$ under the explicitly assumed constant rate $r$ per time unit.
This is a costing convention, not a measured interest model. It includes the
tip and all carried lovelace. A registration's carrier remains locked in the
delivered application output after a successful fold; it is not necessarily
returned to its original funder. A termination can wait without a time bound,
so its capital cost has no finite universal bound here.
For a retraction the requester supplies ordinary lovelace funding for the
ledger fee in addition to the held request value. Its tip is returned, not
earned by a processor. Ownership of the funding is an economic assumption:
the model observes values and addresses, not wallets.

\section{Mechanism and accounting definitions}
\label{sec:mechanism}
\begin{definition}[World and observations]
A world consists of a registry, pending registrations, pending terminations and
a cumulative treasury counter. Its reachability starts from any registry
satisfying the original registry reachability relation, with both pending pools
empty and the counter zero. It then closes under accepted bookings, rejections,
retractions and economic folds. It does not assert that every initial registry
was produced by the economic transactions of this paper.
\end{definition}
An observed economic transaction pairs ledger-shaped inputs, references,
outputs, fee and donation with request contents and typed token movements.
The observation checks bind these fields to an unspent set and the world.
This is a logical observation relation, not a decoder for a Cardano transaction.
The definition is \decl{Singular.WorldReachable}; the accepted judgements are
\decl{Singular.book}, \decl{Singular.bookTermination},
\decl{Singular.rejectionAccepted}, \decl{Singular.retractionAccepted} and
\decl{Singular.foldAccepted}, at the pin above. These are definitions, not
axiom-checked theorem claims; their consequences are cited below.

\begin{definition}[Amounts and clocks]
Let $R$ be the registration charge, $t_{\min}$ the minimum tip, $\delta$ the
maximum booking width, and $W_p,W_r$ the processing and reclaim windows.
These and the registry identity and policy bindings are fixed at creation.
The minimum tip occupies the existing wire key \texttt{maxFee}; it is no
longer an upper fee pin. A request declares tip $t\ge t_{\min}$ and carrier
amount $a$. Its held lovelace is $H$. A registration is booked in a nonempty
interval $[L,U)$ with $U-L\le\delta$, and records $s=U$.
\end{definition}
The carrier amount is declared under the ledger's output-minimum rules; it
is not an additional registry parameter. The model checks exact declared
amounts, not the actual minimum or the fee needed by a concrete transaction.
The original application token's carrier and a later termination request's
carrier need not have equal lovelace amounts: their output shapes differ.

\begin{table}[ht]
\caption{Registration request. Extra assets or hidden lovelace surplus are refused.}
\centering\small
\begin{tabular}{p{.25\linewidth}p{.65\linewidth}}\toprule
Field & Required contents\\\midrule
Lovelace & Exactly $H=R+t+a$; $t\ge t_{\min}$.\\
Payload & Register, key, refund address and destination.\\
Datum & Separately booked destination datum; copied by the fold, not included in these tag-name commitments.\\
Application tag & Exactly one, bound to the payload and nonce; minted in this booking.\\
Trie-ready tag & Exactly one, bound to the operation and the same nonce; minted in this booking.\\
Clock & $s=U$, $L<U$, $U-L\le\delta$.\\
Exits & Fold; signed retraction; reference-read rejection.\\\bottomrule
\end{tabular}
\end{table}
\begin{table}[ht]
\caption{Termination request. Application consent is release from custody.}
\centering\small
\begin{tabular}{p{.25\linewidth}p{.65\linewidth}}\toprule
Field & Required contents\\\midrule
Lovelace & Exactly $H=t+a$; $t\ge t_{\min}$.\\
Singular tag & Exactly the genuine tag of this active key, moved from application custody.\\
Trie-ready tag & Exactly one, bound to its operation and nonce; minted in this booking.\\
Application tag & None.\\
Refund and clock & Neither a refund address nor a booked time.\\
Exit & Fold only, without a registration window.\\\bottomrule
\end{tabular}
\end{table}

\checked The request shapes, same-booking tag minting and bounded booking clock
are consequences of accepted booking and reachable pending membership.
\evidence{booked_request_has_exact_contents}
\evidence{booked_request_tags_commit_to_own_fields}
\evidence{booked_time_is_interval_end}
\evidence{booking_width_within_limit}
\evidence{termination_booking_has_exact_contents}
\evidence{booked_termination_tag_commits_to_own_operation}
Tag names equal the computed commitments; this is not a proof of collision
resistance. Application tags minted outside a booking are outside Singular's
supply guarantees and cannot replace the simultaneous mint required by booking.
The singular tag is a different asset: registration mints it at the destination,
termination booking moves it into a request, and termination folding burns it.
The state token identifies the registry output. The abstract identity spaces
exclude its equality with a trie-ready tag even if a commitment is zero.
\evidence{trie_ready_tag_is_not_state_token}

\subsection{Admitting an exit}
A registration fold's entire nonempty validity interval must lie in
$[s,s+W_p)$. A retraction's must lie in $[s+W_p,s+W_p+W_r)$ and its refund
address must sign. A rejection needs a current lifecycle mismatch, or a
nonempty interval beginning at or after $s+W_p+W_r$. Thus mismatch remains
available during either window, while a compatible request is protected
through both. These are interval-containment rules, not tests of a wall clock
supplied by the processor. Equality of an interval's upper bound with the
window endpoint is allowed because the interval is half-open.

\begin{center}\small
\begin{tabular}{llll}\toprule
Registration time & Fold & Retract & Reject\\\midrule
$[s,s+W_p)$ & Yes & No & Mismatch\\
$[s+W_p,s+W_p+W_r)$ & No & Signed & Mismatch\\
$[s+W_p+W_r,\infty)$ & No & No & Expiry or mismatch\\\bottomrule
\end{tabular}
\end{center}
The table describes admission windows, subject to all other checks. A funded
retraction exists in the model under the explicit premises stated below;
neither a window nor that existence proof guarantees ledger inclusion.

A retraction spends one pending registration and may spend ordinary lovelace
outputs to pay its fee. Extra inputs hold none of the four registry assets,
have neither the request nor state role, and carry no request datum. Spending
inputs are present, distinct and disjoint from references. The state is neither
read nor spent. The refund address signs and receives at least all held
lovelace, including the tip; both request tags burn. No tip is taken and
nothing goes to the treasury. The checked payment bound below states precisely
what ordinary change at another address may contain.

A rejection references the actual unspent state output: one state token with
the exact policy and name, state role, inline datum, complete matching
configuration and current root. It never spends that output. Its distinct
request inputs are exactly its members; every member is checked against that
same state and interval. One invalid member refuses the entire batch. A fold
that consumes the referenced output first invalidates this proposed rejection;
rebuilding can starve indefinitely. A rejected proposal that never enters the
ledger has no ledger transaction fee, though preparation has an external cost.
The state and protection results are stated in Section~\ref{sec:properties}.

The request validator governs spending; a tag policy's permissive burn branch
does not confer a spending right. Registration exits burn both request tags;
a termination fold burns its trie-ready and singular tags. An untagged output
at the request address is not a pending request and can be swept. The checked
sweep judgement preserves the world; a concrete script implementing that
distinction remains a correspondence obligation.
\evidence{sweep_leaves_pool_untouched}

\subsection{Allocation is not net payoff}
\begin{definition}[Fold allocations]
For a fold, let $I,O$ be the sums of ordinary input and output lovelace,
respectively, $f$ its ledger fee, $T=\sum t$ its tips, $n$ its registrations,
and $D$ its donation. State input and continuation values are equal.
Ordinary funding and change are lovelace only; an unrelated native asset is
refused. Carrier outputs hold exactly their declared amount, booked datum and newly
minted singular tag. Ordinary change is a separate output even if its address
equals a carrier's address.
\end{definition}
\checked
\begin{theorem}[Economic fold allocation]
Every accepted economic fold in a reachable world satisfies
\begin{align}
D &= nR+\sum_{q\text{ termination}}(H_q-t_q),\label{eq:donation}\\
O+f &= I+T,\qquad f\le I.\label{eq:gross}
\end{align}
It is balanced; all request tags are burned and the treasury counter increases
by $D$.
\end{theorem}
The mechanized proof is the following checked declaration. The table separates
request allocation from fees.
\evidence{fold_routes_every_amount}

\begin{table}[ht]
\caption{Request allocation. Rejection alone supplies no full balance proof.}
\centering\small
\begin{tabular}{p{.24\linewidth}p{.64\linewidth}}\toprule
Exit & Allocation of held lovelace and token fate\\\midrule
Registration fold & $R$ donation; $a$ carrier at destination; $t$ gross ordinary allocation. Burn both request tags; mint singular tag.\\
Registration rejection & Refund address receives at least $R+a$. The judgement proves no exact processor reward. Burn both request tags.\\
Registration retraction & At least $H$ refunded, including $t$; ordinary funding pays $f$; zero donation, no tip taken. Burn both request tags. Balanced.\\
Termination fold & $H-t=a$ donation; $t$ gross ordinary allocation. Burn trie-ready and singular tags; no refund or destination output.\\\bottomrule
\end{tabular}
\end{table}
Rejection and retraction refund evidence is given under Property~6 below.
The rejection promise is a refund floor, not an exact split. Its judgement
checks a lower bound and does not require full monetary
balance; neither an exact wage nor a rejection analogue of~\eqref{eq:gross}
is supplied. A processor may leave part of its tip to the requester;
over-refunding is compatible with that lower bound.

\checked For a retraction, let $I$ be its ordinary funding, $F$ all output
lovelace at its refund address and $P$ all output lovelace elsewhere. Every
accepted retraction in a reachable world satisfies
\[
 F\ge H,\qquad f\le I,\qquad D=0,\qquad P\le I-f,
\]
and is balanced with an unchanged treasury. Ordinary funding and every output
are lovelace only. The model has no wallet-ownership or output-purpose
observation: $P$ may be change or a payment to a third party, but neither can
consume the request's cost or tip. Saying that the requester pays assumes it
owns the ordinary funding; the checked theorem itself does not assign owners.
\evidence{retraction_pays_only_the_refund_and_the_fee}
\hand Conservation gives $F+P+f=H+I$. Subtracting $F\ge H$ gives the
outside-address bound; any over-refund consumes ordinary funding as well.

\hand For a rejection whose inputs are exactly its members, additionally
assume \decl{Singular.balanced} and zero donation. Partition its outputs into
those paid to the set of refund addresses, totaling $F$, and all other outputs,
totaling $P$. Distinct members and summed floors give $F\ge\sum H-T$.
Balance gives $F+P+f=\sum H$, hence $P+f\le T$, and therefore $P\le T$.
This bounds value outside refund addresses, not the wealth of a party that
also controls a refund address. Balance is an additional ledger fact here,
not a checked conclusion of rejection admission.

\hand With signed differences,~\eqref{eq:gross} gives $O-I=T-f$.
This follows by subtracting $I+f$ in the integers, not by truncated natural
subtraction. It is the net change of the ordinary funding pool. If a processor
owns that pool, pays the fee and has no other role, its monetary net is $T-f$;
including external cost gives $T-f-e$. An address shared with a carrier does
not make the carrier's value a processor reward.

For one request, assume separate funder, destination, processor and treasury;
write $b$ for the booking fee. Measure transfers from before booking, treat
assets and application services separately, and count carrier value as an
asset of the destination. The following illustrative accounting sums to
minus the ledger fees, before external costs. In the funded retraction row,
the funder owns the refund address and ordinary funding, and all remaining
funding returns to it. This is an accepted model construction under the
availability theorem's window, signing and funding premises, not every
possible accepted retraction.
\begin{center}\small
\begin{tabular}{lrrrr}\toprule
Outcome & Funder & Destination & Processor & Treasury\\\midrule
Registration fold & $-R-t-a-b$ & $a$ & $t-f$ & $R$\\
Termination fold & $-t-a-b$ & $0$ & $t-f$ & $a$\\
Rejection, exact refund\! & $-t-b$ & $0$ & $t-f$ & $0$\\
Funded retraction & $-b-f$ & $0$ & $0$ & $0$\\\bottomrule
\end{tabular}
\end{center}
\hand The first two rows sum to $-b-f$ by cancellation. The third does too
\emph{if} an additionally balanced rejection pays exactly $H-t$ to the
funder's refund address, donates zero and returns the residue to its processor.
These added hypotheses are not proved by the rejection judgement. Native
asset values and a preexisting singular tag's valuation are excluded.
The retraction row returns $H+I-f$ to the funder from its total $H+I+b$
outlay, giving $-b-f$; it pays no processor or treasury allocation.
\evidence{pending_registration_can_be_retracted}

For example, with $R=2{,}000{,}000$, $t=400{,}000$, $a=1{,}500{,}000$,
$b=200{,}000$ and $f=170{,}000$, a registration holds $3{,}900{,}000$.
Its four net entries are $-4{,}100{,}000$, $1{,}500{,}000$, $230{,}000$
and $2{,}000{,}000$, summing to $-370{,}000$. These are arithmetic examples,
not measured output minima, costs or recommended parameters.

\subsection{Donation encoding}
The model represents $D$ by a natural number; a zero donation is encoded by
omitting the optional donation field, not by serializing a zero value.
\evidence{zero_donation_omits_field}
Conway transaction-body field 22 is a positive coin donation; optional field
21 is the separate current-treasury-value assertion~\cite{ConwayCDDL}.
When the latter is present, the ledger checks it against current treasury
state~\cite{ConwayLedger}. Donation alone does not make that assertion or
introduce its stale-value hazard. This describes the cited ledger revision;
translation to a concrete transaction still needs correspondence evidence.

\section{Properties at the checked revision}
\label{sec:properties}
Each reference below uses the full pin in Appendix~\ref{sec:artifact} and its
verbatim axiom line. Unless stated otherwise, economic results quantify over
an arbitrary reachable world and an arbitrary transaction accepted by the
relevant judgement, not just a builder's output. The distinction between the
original registry fold and the economic observation is explicit.

\paragraph{1. Fixed parameters.}
\checked An accepted original registry step or batch satisfies
\decl{Singular.onlyRootChanged}: the root may change but all other configuration
fields, including $R,t_{\min},\delta,W_p,W_r$ and policy bindings, are preserved.
\evidence{parameters_preserved_by_fold}
\evidence{parameters_preserved_by_batch}
\hand The economic fold reuses that transition through
\decl{Singular.foldProjectedRegistry} and \decl{Singular.foldBatch}. Inspection of
\decl{Singular.foldAccepted} shows that its resulting registry comes from this
batch. The projection restores carried termination holdings for the original
registry interface and supplies synthetic approval evidence to that interface;
it is not an actual termination application tag or a second mint.
There is no registry update or migration operation in the defined transition
set. This is a statement about that set, not an assertion that an arbitrary
deployed policy cannot spend the state.

\paragraph{2. Rejection preserves the registry.}
\checked Accepted rejection leaves the registry, pending terminations and
singular supply equal. Its observed spending and outputs carry no state or
singular tokens; no such mint is present. Request-tag burns are separate.
The genuine referenced state remains unspent after applying the transaction.
\evidence{rejection_leaves_registry_untouched}
\evidence{rejection_requires_referenced_state}

\paragraph{3. State noninterference, with its actual scope.}
\checked If the \emph{original admitted registry fold} was valid before a
rejection, it remains valid afterward. A supplied list of needed unspent
outputs also remains present if all its other inputs are disjoint from the
rejection's spends. The referenced state is protected by the rejection checks.
\evidence{rejection_cannot_invalidate_fold}
\missing This declaration preserves \decl{Singular.admittedFoldBatch} and
needed outputs. It is not a theorem preserving the whole later
\decl{Singular.foldAccepted} judgement, fees, wallet state or ledger acceptance.
The shared-input counterexample below prevents dropping disjointness.

\paragraph{4. Live protection and a non-backdated clock.}
\checked A pending registration whose key still admits registration cannot
be a member of an accepted rejection whose interval begins before $s+W_p+W_r$.
This includes the reclaim window, not just the processing window. Given an
actual inclusion time $\tau\in[L,U)$, the checked booking rule gives
$\tau\le s$ and $\tau+W_p+W_r\le s+W_p+W_r$.
\evidence{live_request_not_rejected}
\evidence{booked_time_not_before_inclusion}
A current mismatch is outside the compatibility hypothesis: a losing race
can be rejected immediately. The inequality states an expiry threshold, not
an inclusion guarantee or the existence of a retraction transaction.

\paragraph{5. Current evidence is mandatory.}
\checked An accepted rejection has an unspent reference with the exact state
identity, inline configuration and current root; that reference is not spent.
\evidence{rejection_requires_referenced_state}
The per-member mismatch test reads this current state. An alleged booking-root
fault is neither required nor substituted. This rules out evidence-free
acceptance in the economic judgement; it does not verify concrete Merkle code.

\paragraph{6. Refunds are additive per address.}
\checked Rejected members and their input references are distinct. For every
refund address $A$, outputs paid to it total at least
\[
 F_A=\sum_{q:\,\mathrm{refund}(q)=A}(H_q-t_q).
\]
The obligations have one entry per address. The output cannot be counted as a
refund at a different address, and duplicated members cannot reuse it.
\evidence{rejection_batch_members_distinct}
\evidence{rejection_refund_floor}
\evidence{refunds_sum_per_address}
\checked Accepted retraction identifies one pending registration, its signed
refund address and a contained reclaim interval, burns both tags, returns at
least its entire $H$, reads and spends no state output, and removes that
request from the pool.
\evidence{retraction_refund_floor}
\checked The request is among the spending inputs; ordinary funding may
accompany it. The retraction balances, pays its fee from that funding, donates
zero and leaves the treasury unchanged. Value at other addresses is bounded
by ordinary funding minus the fee. Neither the cost nor the tip can pay it.
\evidence{retraction_pays_only_the_refund_and_the_fee}

\paragraph{7. Charge and return.}
\checked An economic fold registering $n$ keys donates at least $nR$;
Equation~\eqref{eq:donation} gives its exact amount. The charge is included in
the rejected address's refund floor and in the retracted request's floor.
\evidence{fold_donates_registration_cost}
\evidence{registration_cost_returned_on_exit}
\hand A successful fold removes its registration request, and the lifecycle
has no operation returning its treasury charge. This concerns the mechanism's
own transfers, not a prohibition on a later treasury payment to the same party.

\paragraph{8. Terminations are unrejectable.}
\checked Every genuine pending termination keeps its key active. Mismatch
therefore fails and there is no expiry ground. It cannot occur as a member of
any accepted rejection, whichever ground is offered.
\evidence{termination_unrejectable}

\paragraph{9. At most one pending termination per key.}
\checked Singular supply counts application custody \emph{plus} actual tags
in pending terminations. It is one exactly for an active key and zero otherwise,
never greater than one. Carried count equals pending termination count, so at
most one is pending for each key.
\evidence{world_supply_sync}
\evidence{at_most_one_pending_termination}
Counting only application custody after release would give the wrong invariant.

\paragraph{10. Termination provenance.}
\checked Every termination selected by an accepted economic fold was pending,
is observed in the unspent set, carries its key's genuine singular tag, and
has an earlier accepted termination booking from a reachable world. That
booking consumed the unique matching custody holding. The application's
release rule itself is an external premise, not rechecked by Singular.
\evidence{only_booked_terminations_fold}
\evidence{termination_request_carries_genuine_token}

\paragraph{11. Conditional model availability.}
\checked For any pending termination in a reachable world, any natural-number
time $u$, funding $I$ and fee $f$ with $f\le I$ and $f\le t$, and a predicate
saying a folder accepts its tip, there exist a coherent unspent set and an
accepted singleton fold with interval $[u,u+1)$. It ends the key as Terminal,
removes the request and produces a reachable world.
\evidence{booked_termination_foldable}
\checked For a pending registration in a reachable world, choose a nonempty
interval wholly inside its booked reclaim window, signers including its
refund address, and ordinary funding $I\ge f$. There exist a coherent unspent
set, a balanced accepted retraction with exactly that interval, signers, fee
and funding total, and a reachable after-world without the request. The full
held lovelace is returned. The theorem permits a positive fee and does not
require it to fit inside the refundable tip.
\evidence{pending_registration_can_be_retracted}
\missing The theorem constructs its own unspent set; it does not find funding
or prove an accepted transaction over an arbitrary live ledger snapshot.
Its witness includes a zero-lovelace state output, illustrating the absent
ledger-minimum premise. It proves no willingness, inclusion or bounded delay.

\paragraph{12. Processor cap is a fold allocation bound.}
\checked For an accepted economic fold, $O\le I+T$, with
$O+f=I+T$ and $f\le I$ as above.
\evidence{processors_take_at_most_the_tips}
\hand Thus net ordinary increase is at most $T$. The formal theorem is about
ordinary outputs and inputs, not ownership, every address receiving funds or
every exit kind. With zero ordinary funding the fee condition forces $f=0$.
No exact rejection reward, positive wage or service incentive follows.

\paragraph{13. A measured minimum tip.}
\proposed Choose $t_{\min}$ to cover the maximum measured single-request exit
cost in a declared envelope, including the reference read and required refund
output. The derivation and assumptions follow in Section~\ref{sec:sizing}.
\missing No checked declaration or completed measurement establishes this
cost cover. Retraction's modeled fee funding is explicit, but a measured
ledger envelope and implementation correspondence remain open. Its refunded
tip is not a processor payment and does not enter a processor-income floor.

\paragraph{Supporting supply statement.}
\checked In every reachable world, the number of trie-ready tags actually
held in the two pending pools equals the total number of pending requests.
Application tags held in pending registrations equal their number, with
exactly one in each. This is not a total application-token supply theorem.
\evidence{tag_supplies_equal_pending}
The world has no separate external trie-ready circulation; lifting its counted
supply to all ledger assets requires mint/spend correspondence. Arbitrary
application minting elsewhere is explicitly outside the counted domain.

\section{Parameter constraints and attack costs}
\label{sec:sizing}
\begin{definition}[Sizing envelope]
An envelope $E$ fixes target ledger fee parameters, transaction and script sizes,
execution budgets, output-minimum rules, supported datum sizes, and assumed
inclusion and response delays. Let $c_x(E)$ bound the measured ledger fee of
a supported singleton exit $x$, and let $e_x(E)$ be a chosen bound on external
processor cost. Neither function is currently measured by this paper.
\end{definition}

\proposed If a processor requires nonnegative net income within $E$, an exit
with gross tip $t$ needs $t-c_x(E)-e_x(E)\ge0$. Consequently a sufficient
conditional floor is
\[
 t_{\min}\ \ge\ \max_{x\in X}\{c_x(E)+e_x(E)\},
\]
where $X$ contains only exits with an implemented, balanced allocation paying
that gross tip. The measured fee-only requirement is the special case $e_x=0$.
The current fold equation supports this arithmetic. Rejection permits a
processor to leave some tip to the requester, so a cost-cover construction
needs to choose its refund and fee allocation explicitly. Thus the formula is
not a completed all-exit implementation result. Retraction returns the tip
and is outside $X$: its checked model construction uses ordinary lovelace
funding $I\ge f$, supplied by the requester in this cost model. Increasing
its refundable tip alone does not pay the fee or a processor. Its ledger fee
still needs measurement in $E$.

\hand If measured fold fee were at most $170{,}000$ and external cost at most
$30{,}000$ lovelace, a $200{,}000$ floor would break even in that hypothetical
envelope. A later fee of $260{,}000$ would give $-90{,}000$ net with the same
external cost. No measurement or recommendation is attached to these numbers.

\proposed The design criterion is $0<\delta<W_p$. For desired booking
inclusion slack $d_b$, processing slack $d_p$ after booked time, and reclaim
response slack $d_r$, require $d_b\le\delta$, $d_p<W_p$ and $d_r<W_r$.
These are conditional feasibility requirements; no consensus inclusion bound
establishes $d_b,d_p,d_r$. The checked booking theorem enforces width at most
$\delta$; it does not establish this entire cross-parameter sizing criterion
for every decoded configuration.

\hand Given $\tau\in[L,U)$ and $s=U$, then $0<s-\tau\le\delta$.
Hence a compatible registration's expiry threshold is between
$\tau+W_p+W_r$ and $\tau+\delta+W_p+W_r$. Adding nonnegative windows to the
checked booking inequalities proves this bound. It is a threshold bound only:
a current mismatch is available earlier, and retraction needs an executable
fee-paying transaction as well as time.

\proposed To make an adversary pay at least a selected treasury charge $C$
for $n>0$ successful junk registrations, require $nR\ge C$, equivalently
$R\ge\lceil C/n\rceil$. This controls only donations of successful folds.
Selecting $C$ is a policy choice, not an optimal price or a universal flood
deterrent. A party assigning value $V(k)>R+t+b+a$ to a key can still choose
to buy it under these simple monetary assumptions.

\hand The following costs follow only under the expressly stated allocation
and ownership assumptions; none is a game-theoretic equilibrium claim.
Successful junk folding pays $R$ per key to the treasury, plus ledger fees;
if the funder and processor coincide, the tip is an internal transfer.
If that same party also controls the destination, its carrier remains an
asset, with an opportunity cost rather than an immediate wealth loss.
For a balanced exact-refund self-rejection, the tip is likewise internal;
the combined party loses the booking and rejection fees, plus external and
capital costs. There is no positive treasury floor for this attack.
For a losing race processed by someone else under the same exact-refund
assumption, the funder loses $t+b$ while receiving $R+a$ back. A larger tip
does not buy priority. A termination fold irrevocably donates its declared
carrier $a$; with no guaranteed delay bound, $rH\Delta$ remains unbounded.
For the funded retraction construction with all change returned to its
funder, that party loses $b+f$ plus external and capital costs; it recovers
the registration cost, carrier and tip. There is no treasury charge on that exit.

No bound relates the number of harmed requesters to the number of charged
attacker folds. A griefing ratio would divide total victim losses, including
delay costs, by total attacker expenditure under an explicit schedule. Neither
that schedule nor a finite delay bound is supplied; no finite universal ratio
is derived from a maximum batch size.

Policy identities and the registry identity are separation parameters, not
economic prices. They must identify the intended application and asset policies;
choosing them incorrectly cannot be repaired by a larger tip or window.
Carrier values must satisfy actual output minima when outputs are created.
They are not sized by a fixed registry constant. Changes outside $E$ require
new measurements and may motivate a new registry; no automatic migration or
rescue of old pending requests is implied.


\section{Related work}
The extended UTxO work supplies the state-machine and invariant setting in which
this construction lives~\cite{ExtendedUTXO20}. Singular specializes it to one
permanent lifecycle, authenticated request admission and explicit token custody.
It does not propose a new ledger model or prove the ledger implementation.

Roughgarden studies transaction-fee mechanisms, including the incentives of
users and block producers~\cite{Roughgarden21}. Singular's minimum tip is instead
an immutable application-level floor. Its accounting equation is neither an
incentive-compatibility theorem nor a block-allocation mechanism. A processor
may prefer another transaction even when this one covers its immediate fee.

The transaction-ordering attacks documented in Flash Boys 2.0 concern value
obtained by ordering and front-running~\cite{FlashBoys19}. Singular protects a
specific refund floor after a lost race; it neither hides the request nor
bounds the value of capturing its key. An application's authorization policy
must provide any ownership restriction.

Namecoin's empirical namespace study exposes the difficulty of useful allocation
and squatting~\cite{Namecoin15}. ENS separates a registry, resolvers and
registrars, leaving allocation policy to registrars~\cite{EIP137}. Singular
likewise leaves application authorization external, but its registry leaf is a
permanent lifecycle marker and its representative is an asset held in application
custody. The registration charge is not a rent, auction or proof of a right to a
human-readable name. It establishes no optimal namespace allocation.

Lean provides the language and proof kernel used to check the statements
in this paper~\cite{Lean4}. The contribution is a reviewable separation of lifecycle,
request custody, reference-read rejection and conditional accounting, together
with revision-bound proofs and explicit missing correspondence. A proof of this
model is not a proof of its adequacy for every application.

\section{Limitations and implementation standing}
The validators and builders are \emph{not yet conformant} to the ruled economics.
A model result is not a hosted conformance result, an accepted transaction,
a deployed validator identity or an application's integration acceptance.
The retained implementation's agreement with its older driver does not close
these requirements. Coverage must come from observed receipts, with unmet
requirements kept visible.

The application's release rule, custody integration, and adaptations of
open-datum and cardano-keri remain open. Open-datum's earlier control-signature
and deposit-return behavior needs reconciliation with release at booking and
termination donation. Cardano-keri's reaping, refund and convictor reward need
their own correspondence design. Singular supplies none of those application
promises merely by recognizing an asset policy.

Measured registration cost, minimum tip, booking-width bound and their target
ledger envelope remain open. No carrier amount is a registry parameter:
the requester declares it under ledger minimum rules. Parameter sufficiency
need not survive changes in fees, output sizes, minimums or exchange rates.
No selected number is claimed optimal.

Ordinary funding and change are lovelace only: both the economic fold and
retraction refuse an unrelated native asset. Retraction has checked balance
and conditional availability, but neither statement supplies wallet funds,
computes a ledger minimum, proves builder success or guarantees inclusion.
The model cannot distinguish change at another address from a third-party
payment. Such outputs can consume only ordinary funding left after the fee,
never the request's registration cost, carrier or tip.

\subsection{Counterexamples to stronger readings}
\hand \textbf{No universal flood price.} Book many registrations without
folding them. No $R$ donation occurs. Even eventual balanced self-rejection
can make the tip internal, leaving only fees and opportunity cost. Successful
junk registrations do pay $R$ but still consume state-output turns.

\hand \textbf{Protection is not priority.} Two authorized registrations for
the same unknown key can be pending. The first fold makes it active; the
second is now mismatching and may be rejected before its time expiry. Its
refund floor survives, not a right to the key or the full tip.

\hand \textbf{Noninterference needs disjoint inputs.} If a proposed fold and
a rejection spend the same pending registration, rejection removes that
input. Preserving the state output alone cannot preserve the whole fold.
Conversely repeated folds can invalidate a rejection's state reference before
it lands, without violating the rejection theorem.

\hand \textbf{No positive wage.} Choose fold tip $t=100{,}000$, fee
$f=170{,}000$ and ordinary funding $I\ge170{,}000$. The allocation equation
gives net ordinary change $-70{,}000$. Fold acceptance requires funding but
does not impose the availability theorem's extra $f\le t$ hypothesis.

\hand \textbf{No ownership or squatting defence.} An application that
authorizes every request also authorizes a stranger's registration of a
desirable key. Paying $R$ does not prove a right to its name. A tag's unique
custody is an asset invariant, not an ownership adjudication.

\missing \textbf{Full refund is not fee-paying retractability.} A single
$3{,}000{,}000$ input cannot return that full amount and also pay $170{,}000$.
The current accepted-retraction theorem and this conservation counterexample
coexist because its judgement omits balance and excludes a funding input.
Changing the model or claiming real-ledger retraction requires a separate
ruling and new evidence; the floor is not silently weakened here.


\section{Future work: reusable names and durable evidence}
\proposed The next lifecycle extension would expose the five other edges.
They are refused by the active contract and are not proved functionality
here. One group is non-terminality: deleting an active key would let an
application release a name for reuse. That would be a different lifecycle
contract, not an exception to permanence in an existing registry.

\proposed The second group aims at proof durability through witnesses.
An absent witness would attest that a name is not taken; the proposed path
also includes changing that absent state to active and deleting the absent
state. A terminal witness would be minted on demand and attest a termination
permanently. These four edges together with deletion of an active key account
for the five proposed extensions.

\proposed A separate presence witness is also proposed. It would be minted
whenever the registry mints a name's token, and burned when that name is
terminated or deleted. Thus the three proposed attestations concern absence,
presence and termination. None is supplied as a working witness mechanism
by this paper's active two-operation model.

\hand The motivation is the lifetime of a referenced output. A proof used
by a transaction against an application output holding the name's token
depends on that output remaining unspent; an application update can replace
it. A proof against the registry's current root likewise depends on the
state output, which a fold replaces. Once that referenced output is spent,
the transaction cannot use it. This explains the desire for more durable
attestations, without proving that a proposed token protocol achieves it.
\missing The boundary model now specifies the three policy decisions and
terminal discard. Issuance through the proposed lifecycle extensions,
invalidation, custody and concrete verification still require model and
implementation evidence. A presence or absence witness cannot be promised valid after
its own condition has ceased to hold.

\clearpage
\appendix
\section{Reproducible model evidence and trusted base}
\label{sec:artifact}
All entries below use revision
\begin{center}\small\texttt{\modelpin}\end{center}
with Lean 4.27.0. The command \texttt{lake env lean tools/axioms.lean}
was run over an archive of that revision using the repository's seeded compiled
modules at the same model source. Every listed declaration was separately
confirmed by \texttt{git show} at the pin. The raw report's SHA-256 is
\begin{center}\scriptsize\nolinkurl{1f8cf8a5c792a9419da90ef0ab465923248a74914790985c10943b71c1621a08}\end{center}
No listed result uses \texttt{sorryAx} or an additional nonstandard axiom.
The following are verbatim axiom lines, with typographic wrapping only.
Source paths are relative to the repository at the stated revision.
\begingroup\footnotesize
\begin{longtable}{p{.94\linewidth}}
\toprule Declaration evidence at the model pin\\\midrule\endhead
\texttt{AXIOMS }\nolinkurl{Singular.Statements.at_most_one_pending_termination}\texttt{ [propext, Quot.sound, Classical.choice]}\\
\nolinkurl{lean/Singular/Statements/Termination.lean:141}\\[5pt]
\texttt{AXIOMS }\nolinkurl{Singular.Statements.booked_request_has_exact_contents}\texttt{ [propext, Quot.sound, Classical.choice]}\\
\nolinkurl{lean/Singular/Statements/Booking.lean:13}\\[5pt]
\texttt{AXIOMS }\nolinkurl{Singular.Statements.booked_request_tags_commit_to_own_fields}\texttt{ [propext, Quot.sound, Classical.choice]}\\
\nolinkurl{lean/Singular/Statements/Booking.lean:192}\\[5pt]
\texttt{AXIOMS }\nolinkurl{Singular.Statements.booked_termination_foldable}\texttt{ [propext, Quot.sound, Classical.choice]}\\
\nolinkurl{lean/Singular/Statements/FoldTermination.lean:36}\\[5pt]
\texttt{AXIOMS }\nolinkurl{Singular.Statements.booked_termination_tag_commits_to_own_operation}\texttt{ [propext, Quot.sound, Classical.choice]}\\
\nolinkurl{lean/Singular/Statements/Termination.lean:221}\\[5pt]
\texttt{AXIOMS }\nolinkurl{Singular.Statements.booked_time_is_interval_end}\texttt{ [propext, Quot.sound, Classical.choice]}\\
\nolinkurl{lean/Singular/Statements/Booking.lean:82}\\[5pt]
\texttt{AXIOMS }\nolinkurl{Singular.Statements.booked_time_not_before_inclusion}\texttt{ [propext, Quot.sound, Classical.choice]}\\
\nolinkurl{lean/Singular/Statements/Booking.lean:127}\\[5pt]
\texttt{AXIOMS }\nolinkurl{Singular.Statements.booking_width_within_limit}\texttt{ [propext, Quot.sound, Classical.choice]}\\
\nolinkurl{lean/Singular/Statements/Booking.lean:140}\\[5pt]
\texttt{AXIOMS }\nolinkurl{Singular.Statements.fold_donates_registration_cost}\texttt{ [propext, Quot.sound, Classical.choice]}\\
\nolinkurl{lean/Singular/Statements/FoldAccounting.lean:11}\\[5pt]
\texttt{AXIOMS }\nolinkurl{Singular.Statements.fold_routes_every_amount}\texttt{ [propext, Quot.sound, Classical.choice]}\\
\nolinkurl{lean/Singular/Statements/FoldAccounting.lean:25}\\[5pt]
\texttt{AXIOMS }\nolinkurl{Singular.Statements.live_request_not_rejected}\texttt{ [propext, Quot.sound, Classical.choice]}\\
\nolinkurl{lean/Singular/Statements/RejectionProtection.lean:13}\\[5pt]
\texttt{AXIOMS }\nolinkurl{Singular.Statements.only_booked_terminations_fold}\texttt{ [propext, Quot.sound, Classical.choice]}\\
\nolinkurl{lean/Singular/Statements/FoldTermination.lean:13}\\[5pt]
\texttt{AXIOMS }\nolinkurl{Singular.Statements.parameters_preserved_by_batch}\texttt{ [propext, Quot.sound, Classical.choice]}\\
\nolinkurl{lean/Singular/Statements/Preservation.lean:19}\\[5pt]
\texttt{AXIOMS }\nolinkurl{Singular.Statements.parameters_preserved_by_fold}\texttt{ [propext, Quot.sound, Classical.choice]}\\
\nolinkurl{lean/Singular/Statements/Preservation.lean:11}\\[5pt]
\texttt{AXIOMS }\nolinkurl{Singular.Statements.processors_take_at_most_the_tips}\texttt{ [propext, Quot.sound, Classical.choice]}\\
\nolinkurl{lean/Singular/Statements/FoldAccounting.lean:50}\\[5pt]
\texttt{AXIOMS }\nolinkurl{Singular.Statements.refunds_sum_per_address}\texttt{ [propext, Quot.sound, Classical.choice]}\\
\nolinkurl{lean/Singular/Statements/Refunds.lean:92}\\[5pt]
\texttt{AXIOMS }\nolinkurl{Singular.Statements.registration_cost_returned_on_exit}\texttt{ [propext, Quot.sound, Classical.choice]}\\
\nolinkurl{lean/Singular/Statements/Refunds.lean:70}\\[5pt]
\texttt{AXIOMS }\nolinkurl{Singular.Statements.rejection_batch_members_distinct}\texttt{ [propext, Quot.sound, Classical.choice]}\\
\nolinkurl{lean/Singular/Statements/RejectionState.lean:91}\\[5pt]
\texttt{AXIOMS }\nolinkurl{Singular.Statements.rejection_cannot_invalidate_fold}\texttt{ [propext, Quot.sound, Classical.choice]}\\
\nolinkurl{lean/Singular/Statements/RejectionState.lean:67}\\[5pt]
\texttt{AXIOMS }\nolinkurl{Singular.Statements.rejection_leaves_registry_untouched}\texttt{ [propext, Quot.sound, Classical.choice]}\\
\nolinkurl{lean/Singular/Statements/RejectionState.lean:12}\\[5pt]
\texttt{AXIOMS }\nolinkurl{Singular.Statements.rejection_refund_floor}\texttt{ [propext, Quot.sound, Classical.choice]}\\
\nolinkurl{lean/Singular/Statements/Refunds.lean:12}\\[5pt]
\texttt{AXIOMS }\nolinkurl{Singular.Statements.rejection_requires_referenced_state}\texttt{ [propext, Quot.sound, Classical.choice]}\\
\nolinkurl{lean/Singular/Statements/RejectionState.lean:43}\\[5pt]
\texttt{AXIOMS }\nolinkurl{Singular.Statements.retraction_refund_floor}\texttt{ [propext, Quot.sound, Classical.choice]}\\
\nolinkurl{lean/Singular/Statements/Refunds.lean:30}\\[5pt]
\texttt{AXIOMS }\nolinkurl{Singular.Statements.sweep_leaves_pool_untouched}\texttt{ [propext, Quot.sound, Classical.choice]}\\
\nolinkurl{lean/Singular/Statements/RejectionState.lean:105}\\[5pt]
\texttt{AXIOMS }\nolinkurl{Singular.Statements.tag_supplies_equal_pending}\texttt{ [propext, Quot.sound, Classical.choice]}\\
\nolinkurl{lean/Singular/Statements/TagSupply.lean:13}\\[5pt]
\texttt{AXIOMS }\nolinkurl{Singular.Statements.termination}\texttt{ [propext, Quot.sound, Classical.choice]}\\
\nolinkurl{lean/Singular/Statements/Preservation.lean:223}\\[5pt]
\texttt{AXIOMS }\nolinkurl{Singular.Statements.termination_booking_has_exact_contents}\texttt{ [propext, Quot.sound, Classical.choice]}\\
\nolinkurl{lean/Singular/Statements/Termination.lean:13}\\[5pt]
\texttt{AXIOMS }\nolinkurl{Singular.Statements.termination_request_carries_genuine_token}\texttt{ [propext, Quot.sound, Classical.choice]}\\
\nolinkurl{lean/Singular/Statements/Termination.lean:151}\\[5pt]
\texttt{AXIOMS }\nolinkurl{Singular.Statements.termination_unrejectable}\texttt{ [propext, Quot.sound, Classical.choice]}\\
\nolinkurl{lean/Singular/Statements/RejectionProtection.lean:52}\\[5pt]
\texttt{AXIOMS }\nolinkurl{Singular.Statements.trie_ready_tag_is_not_state_token}\texttt{ [propext]}\\
\nolinkurl{lean/Singular/Statements/Booking.lean:185}\\[5pt]
\texttt{AXIOMS }\nolinkurl{Singular.Statements.world_supply_sync}\texttt{ [propext, Quot.sound, Classical.choice]}\\
\nolinkurl{lean/Singular/Statements/Termination.lean:54}\\[5pt]
\texttt{AXIOMS }\nolinkurl{Singular.Statements.zero_donation_omits_field}\texttt{ [propext, Quot.sound, Classical.choice]}\\
\nolinkurl{lean/Singular/Statements/Accounting.lean:82}\\[5pt]
\bottomrule
\end{longtable}
\endgroup


The trusted base includes the Lean kernel and the printed standard axioms,
the adequacy of the mathematical specification, and the correspondence of
observed values, identities and validity intervals to actual ledger data.
The executable map and commitment functions are stand-ins, not cryptographic
implementations. In particular, economic tag accounting counts repeated names
and collisions and does not assume injective tag hashes; this does not establish
cryptographic binding of a real policy. The logical current-state read is not
a proof of concrete Merkle-path verification. The legacy model's descriptive
collision-free profile must not be mistaken for a proof of its executable hash.

Application authorization and custody release are supplied evidence boundaries.
The Lean-to-validator correspondence is an assumption requiring conformance
work, not a result of these proofs. Ledger minimums, fee measurement, inclusion
and application integration remain outside the checked economic statements.

\clearpage
\bibliographystyle{splncs04}
\bibliography{refs}
\end{document}
